\section*{Layman Summary}

%MARTIN GOOD
% Many products are made using very tiny pieces of material called nanoparticles. These are found in things like plastics, medicines, and food. To make sure these products are safe and good quality, we need to measure what these tiny pieces are like. For example, their size and how they behave.

% However, nanoparticles are so small that we cannot see them clearly, even with a microscope that uses light. So we need to find new ways to study them.

% This thesis develops an alternative method to do this. It measures where a particle is located in space, and it watches how that position changes over time. Using mathematics, we can then learn more about the particle from the way it moves. This helps us measure nanoparticles more accurately.

% CRINGE
The thermal motion of particles has been extensively studied to understand the influence of temperature on particle interactions. Thermal behavior is typically characterized by the temperature of a body or environment, which represents the average kinetic energy of a large ensemble of particles and follows the Maxwell–Boltzmann statistical distribution. However, the influence of temperature on the behavior of individual particles within a complex environment remains less understood.

A more detailed understanding of thermal motion at the single-particle level requires the ability to observe and analyze the trajectories of individual particles. Single-particle tracking (SPT) provides this capability through particle labeling, high-resolution imaging, and computational trajectory reconstruction, enabling direct measurement of particle dynamics with high spatial and temporal resolution.

This thesis presents a method for three-dimensional single-particle tracking using interferometric microscopy. The proposed approach enables precise localization of individual particles in three dimensions and reconstructs their trajectories over time.

Analysis of the reconstructed trajectories allows the extraction of important physical parameters, including the mean squared displacement (MSD), diffusion coefficient, drift velocity, interaction forces, particle size, and confinement effects. These measurements provide quantitative insight into the thermal motion and interactions of individual particles, offering a powerful tool for investigating nanoscale processes that cannot be resolved through ensemble-averaged measurements alone.


%better
% In scientific manufacturing, such as plastics, pharmacology, and food science small molecules, known as nano-particles, are part of the product or procedure. In order to ensure that the output product is of high-quality or up to standard, one tends to want to measure the properties of these nanoparticles. However, due to their size being smaller than that resolvable by light alone, there is a need to develop new methods to do this. This thesis aims to develop one such method to measure three-dimensional particle position in space, and mathematically extract information from this position in time to better improve measurements of nanoparticle characteristics.

% intro
%Particle confinement allows for extended measurements of suspended biomolecules and nanoparticles using optical microscopy as it keeps it within the field of view.

\newpage
\section*{Abstract}
The understanding of thermal movement of paricles within a confined space has been studied extensively in order to understand biological, chemical, and physical processes. Although the theoretical understanding of this confined diffusion is well defined in literature, not much work has been done to understand the effects of electrostatic and hydrodynamic contributions to the size measurements measured from the motion of nanoparticles, as measurements are often intrinsically performed under confinement. This is due to the fact that 3D particle-tracking is not as commonly used as other, higher throughput sizing techniques. By employing developed techniques of nanoparticle tracking in tandem with a full theoretical model, the non-trivial effects of confinement on size-diffusivity were directly investigated and characterized. The work described in the thesis advances the method and tools of 3D nanoparticle analysis both theoretically and in application.

% This thesis provides an understanding and investigation into confined diffusion measurements.

The first section of the thesis describes a theoretical model of Brownian motion based on the confined-Langevin model using a Debye-Hückel electrostatic potential. The model was implemented, based on literature \cite{millanNumericalSimulationsConfined2023}, into a computer-based Python simulation in order to parameterize and create prior distributions of particle motion. 

The second section describes a modified method of 3D tracking to obtain faster calculations of particle position from optical microscopy images. This was based on an algorithm described in literature \cite{dewitPreciseCharacterizationNanometerscale2023}. The validation of the method was done on exact point-spread function simulation data with an accuracy of 1 nm.  

The third section describes the acquisition of  experimental data using a custom interferometric scattering microscope (iSCAT) with image pixel resolution 55 nm acquiring images between 400 and 1000 frames per second. The data consisted of images of scattering fringes from nanoparticle species and was analyzed using the algorithm developed in section 2. The resultant information was accurate 3D particle tracks. The models developed in section 1 were applied to the tracks to obtain infer particle size and charge. 

The thesis presents a method of finding particle size, motion, and electrostatics which can present additional tools to analyze nanoparticle characterization, which was the goal of the thesis. 

\section*{Acknowledgements}

I would like to thank for the guidance, help and resources from Dr. Sabrina R. Leslie, who is also my supervisor of thesis. In addition, I would like to thank Dr. Erik Olsén, a post-doc in the Leslie-Lab, who collaborated and provided guidance to me during this project.